\subsection{Tool corpus}
We snapshot the official MCP registry~\cite{mcpregistry} on 5~October~2026
(\NumRegistry{} latest-version server entries). \NumCandidates{} entries advertise a
streamable-HTTP endpoint that requires no API-key header. For each we send only the
protocol handshake and \texttt{tools/list}; no tool is invoked. \NumRespondingServers{} servers
answer, returning \NumRawTools{} tool definitions (of the rest, \NumFailAuth{} refuse with HTTP 401
and \NumFailRate{} are rate-limited). We drop servers whose tool set is identical to one already kept, and
cap every publisher namespace at 25 servers: two publishers alone had generated thousands of
near-identical wrapper servers. The resulting catalog has \NumCatTools{} tools from
\NumCatServers{} servers and \NumPublishers{} publishers.

\subsection{Operation profiles and sibling groups}
We assign each tool a side-effect severity on a four-level, study-specific scale:
\emph{read}~(0) returns information only; \emph{create}~(1) adds state or sends something;
\emph{update}~(2) mutates existing state, executes, pays or starts/stops; \emph{destructive}~(3)
deletes, revokes or overwrites irreversibly. The scale orders consequences for a typical request;
we do not claim every create is safer than every update.

\emph{Rule profile.} The tool name is split into tokens; the first unambiguous verb from a lexicon
of about 400 verbs sets the class, with noun-like verbs (\texttt{issue}, \texttt{post},
\texttt{order}) accepted only in leading position and the description's first verb as fallback.
The resource is the first noun after the verb, singularised and alias-normalised
(\texttt{repo}$\rightarrow$\texttt{repository}).
\emph{System profile.} A decisive annotation (\texttt{readOnlyHint=true} or
\texttt{destructiveHint=true}) overrides the rule; otherwise the rule stands. It costs no model
call.
A \emph{sibling group} is a set of at least two tools on one server sharing a resource and
covering at least two severities; it is \emph{escalation-capable} if it contains a read tool and
an update or destructive tool. Such groups often also hold create tools, and read-to-write counts a
create as a write.

\emph{Validation.} An AI assistant working under the first author's direction, not one of the evaluated
models, labeled the true operation of a seeded sample of \NumAuditTools{} tools stratified by rule
class, and we compare profiles against those labels (Table~\ref{tab:profiles}). On \RaterToolN{} of
these tools the first author's blind labels agree with the assistant's on \RaterToolAgree\,\%
($\kappa=\RaterToolKappa$), so Table~\ref{tab:profiles} compares profilers with one another more
reliably than it measures their absolute accuracy.

\subsection{Requests}
From the escalation-capable groups we draw 700 with at most three per server and six per
publisher. For each group we target its read tool and up to two of its stronger tools. A
generator (Qwen3.5-9B) sees the target and its siblings and writes one request that only the
target fulfils, without naming it; requests containing a tool name or identifier are discarded.
A validator from a different model family (Gemma3-12B) then routes each request among the shuffled
siblings; we keep a request only if it picks the target. This yields \NumQueries{} requests
(\NumQueriesRead{} generated for a read target, \NumQueriesWrite{} for create, update or destructive
targets) over \NumQueryServers{} servers. By the reference severity of the target, \ReadReqs{} are
reads; read-to-write uses this count. Validating within the sibling group makes each request unambiguous
\emph{among the tools that matter}, so measured escalations are not artefacts of vague requests.

\subsection{Reference severities}
Escalation is scored with a reference severity independent of the system profile: Gemma3-12B
classifies the tool from its name and description, and a decisive annotation overrides it.
Against these LLM-assisted labels the reference reaches \RefExact{}\,\% exact agreement
($\kappa=\RefKappa$).

\subsection{Catalogs, retrievers and selectors}
Experiments run on a \NumPoolTools{}-tool study pool: every tool of the
\NumQueryServers{} servers behind the requests, plus seeded random servers. A per-request catalog
of size $N$ always installs the target's whole server, as a user who connected that server would,
then adds other pool servers in seeded random order. We vary $N$ over 200, 1,000, 5,000 and \NumPoolTools{} tools; for the 20 requests whose server alone exceeds 200 tools (4 above 1,000), the
catalog is that server.
Each tool is indexed as its name, description and parameter names. We compare four retrievers,
BM25, MiniLM-L6~\cite{sbert}, BGE-base~\cite{bge} and Qwen3-Embedding-0.6B~\cite{qwen3emb}, with
exact cosine search for the dense ones; the top $k{=}5$ go to the selector. Selectors are Qwen3.5-9B, Gemma3-12B, Phi-4-14B and Qwen3-1.7B, four open models from three
families (Q4, Ollama, temperature 0), prompted to return one tool name or \texttt{NONE}.

\subsection{Selection variants}
\begin{itemize}
\item \textbf{plain}: top-$k$ names and descriptions.
\item \textbf{native}: plain plus each server's own annotation, when present
(\texttt{[READ-ONLY]}, \texttt{[DESTRUCTIVE]}, \texttt{[WRITE]}).
\item \textbf{tags}: plain plus the system profile's tag on every candidate, and an instruction
not to exceed the requested side effect.
\item \textbf{sibling}: tags, plus \emph{sibling expansion} (each candidate's sibling group is
added, up to $k{+}10$ tools) and \emph{intent ordering}: an intent profiler, the same LLM asked
for \texttt{read|create|update|delete}, puts candidates whose severity matches first.
\item \textbf{expand} and \textbf{expand+tags} (ablations, Qwen3.5 only): sibling expansion in plain
similarity order, without and with tags, and without intent ordering.
\end{itemize}
\textbf{Escalation guard.} After any variant, if the chosen tool's system severity exceeds the
intent's severity by at least $\tau$, the selection is held for user confirmation.
We compare five policies: any increase, an increase of two or more levels, a read-guard (read
requests may not run update or destructive tools), a destructive-guard, and the union of the last
two. The intent is inferred by the selector's own LLM (Qwen3.5 for retriever-only routing), asked
for \texttt{read|create|update|delete}; requests whose intent cannot be parsed (under 2.5\,\%) are
never held.

\subsection{Metrics}
Let $s(\cdot)$ be the reference severity, $g$ the target and $p$ the selected tool. All metrics are
\emph{end to end}: a request whose target the retriever never shortlists counts as wrong even though
the selector never saw the target (Section~\ref{sec:split} separates the stages). We measure
\emph{selection-time potential} escalation: the tool chosen, not executed.
\textbf{Read-to-write rate}, the primary safety outcome: the share of read requests whose selected
tool is a verified escalation, that is, would create, send, overwrite, start or delete something;
the \textbf{side-effect escalation rate} \seer $=\Pr[p\neq g \wedge s(p)>s(g)]$, the broader secondary
metric, computed automatically with the reference or over verified escalations only;
\textbf{exact accuracy} $=\Pr[p=g]$, and \textbf{coarse operation-equivalence accuracy}, which also
accepts a tool on another server with the same operation, resource and severity (diagnostic only:
tenant, scope and parameters can still differ);
\textbf{unblocked \seer} counts selected escalations the guard would let through;
\textbf{hold rate} and \textbf{needless holds} ($p=g$ but held) measure the guard's cost;
\textbf{Exposure@k} is, for read requests, the share whose top-$k$ shortlist contains an update or
destructive tool.

\textbf{Human labels.} The reference and the system profile both err, so a flag counts only when a
human confirms that the selected tool would create, send, overwrite, start or delete more than the
request asked for; we call such escalations \emph{verified}. The first author labeled every flagged (request, tool) pair in the runs we report (all top-5 and
candidate-budget selection runs at 1,000 and \NumPoolTools{} tools, and Qwen3-Embedding top-1 at
all four catalog sizes),
blind to model and variant. The second author, who took no part in building the benchmark or the
first labeling pass, served as independent rater: blind to source, model, variant and the first
author's labels, this rater labeled a seeded sample of 55 items: 25 flagged pairs and 30 unflagged wrong picks,
15 from LLM selection and 15 from retriever-only routing, at most one per request. Both followed
the same written rubric: a pick escalates if, beyond what was asked, it would create a record;
send, submit or post something; overwrite or change a value; start a workflow, job, payment or
purchase; or delete or cancel something, judged from the tool's description, not from a
server's \texttt{destructiveHint} alone. After agreement was computed, the two reviewed every disagreement together
against the full tool descriptions; a label changed only if both agreed, otherwise the case is
\emph{ambiguous}. The consensus is the final label on the double-rated items.

\textbf{Sensitivity estimate.} The reference can also \emph{miss} escalations (a wrong pick whose
reference severity is not higher). We estimate the total rate as verified \seer plus the missed share
times the number of unflagged wrong picks, separately for LLM selection and retriever-only routing,
with the missed share from the consensus sample (Jeffreys estimate, so a zero count still yields a
range). Its 95\,\% range comes from a two-stage bootstrap: servers are resampled as clusters, and the
missed share is redrawn from its posterior. Selector estimates use plain-prompt runs and the retriever estimate its top-1, where the
samples were drawn. Other intervals are 95\,\% Wilson; accuracy differences between variants use paired exact
McNemar tests.
